\section{The comparison with Galois and the gentle group}
The factorization also makes the nonabelian equality question and the
role of gentle characters precise.

\begin{proposition}\label{prop:normalized-comparison}
There is a natural homeomorphism of right-coset spaces
\[
 N/K\longrightarrow G/H,\qquad nK\longmapsto nH.
\]
In particular $G=H$ if and only if $N=K$.
Put $N_0=\ker(\lambda:\GTz\to U)\cap\ker\rho_2$.
Then $\GTz=N_0H$. If a continuous homomorphism from $\GTz$
to any topological group is trivial on $H$, its restriction to $G$
is trivial if and only if its restriction to $N$ is trivial.
\end{proposition}
\begin{proof}
Surjectivity of the coset map is $G=NH$. Two elements of $N$ have
the same right $H$-coset precisely when their quotient lies in
$N\cap H=K$. The map is continuous; compactness of the source and
the Hausdorff property of the target make it a homeomorphism.
There is a single coset exactly when $N=K$, proving the equivalence.

For $g\in\GTz$, choose $h\in H$ with $\lambda(h)=\lambda(g)$.
By \cite[Lemma~7(c)]{H17}, the value $\rho_2(gh^{-1})$ is even;
that lemma uses only the three-cycle equation. The Galois Kummer
image over $\Q^{\mathrm{ab}}$ is $2\Zh$, as proved in
Proposition~\ref{prop:arithmetic-affine}. Choose
$k\in G_{\Q^{\mathrm{ab}}}$ with
$\rho_2(k)=\rho_2(gh^{-1})$. The cocycle law gives
$gh^{-1}k^{-1}\in N_0$, proving the gentle factorization.
Finally $G=NH$ shows directly that a homomorphism killing $H$
kills $G$ exactly when it kills $N$.
\end{proof}

The equivalence is about equality of the two full profinite groups.
The quotient $B$ records its abelian consequence. In particular a
gentle character trivial on Galois proves strictness of $G\supset H$
only if it is nonzero on $N$; a character vanishing on the pentagon
group does not provide such a witness.
